Our goal in this Section is to define quotients of Boolean inverse monoids. In general, an \emph{ideal}\index{Ideal (of a semigroup)} of an inverse semigroup $S$ is a subset $I\subseteq S$ which satisfies $SI\cup IS\subseteq I$. The quotient $S/I$ could be defined as the quotient semigroup of $S$ by the smallest congruence which makes all elements of $I$ equivalent, in which case we obtain the \emph{Rees factor} of $S$ by $I$ (\cite{MR0002893}) and the congruence is given by collapsing all elements of $I$ into a zero element, and nothing else:
\[s\sim t\iff s=t\text{ or }\left\{s,t\right\}\subseteq I\]
However, the Rees factor does not preserve the order-theoretic operations of inverse semigroups.

\begin{example}
Let $X=\left\{a,b\right\}$ be a set with two elements and $S=\mathcal{P}(X)$ the lattice of subsets of $X$. Consider the ideal $I=\left\{\left\{a\right\},\varnothing\right\}$, and let $T$ be the Rees factor of $S$ by $I$ and $\pi:S\to T$ the quotient map.

Then $\left\{a\right\}\lor\left\{b\right\}=X$, however $\pi(\left\{a\right\})\lor\pi(\left\{b\right\})=\pi(\left\{b\right\})$, which is not $\pi(X)$.
\end{example}

Therefore we need to consider more specific classes of ideals and a more suitable congruence. We will be interested in \emph{$\lor$-ideals}, as defined in \cite{MR3448402}. (These are called $\lor$-\emph{closed} ideals in \cite{MR3628219,MR3519501,MR3715848}, \emph{tightly closed} in \cite{MR2974110} and \emph{additive} in \cite{MR3700423}).

\begin{definition}\index{Ideal (of a semigroup)!$\lor$-ideal}%\index{Ideal}
A \emph{$\lor$-ideal} (read \emph{join ideal}) of a distributive inverse semigroup $S$ is a nonempty subset $I\subseteq S$ satisfying
\begin{enumerate}
\item[(i)] For all $s\in S$ and $i\in I$, $si\in I$ and $is\in I$ (i.e., $I$ is \emph{absorbing});
\item[(ii)] For any two compatible elements $i,j\in I$, we have $i\lor j\in I$.
\end{enumerate}
\end{definition}

Quotients of (weakly) Boolean inverse semigroups, under ``additive congruences'', are described in \cite[Section 3.4]{MR3700423}, and in particular their quotients by $\lor$-ideals in \cite[Proposition 3.4.6]{MR3700423}. Although this is the case of most interest, and in which most proofs can be significantly simplified, we initially consider the larger class of distributive inverse semigroups.

Given a $\lor$-ideal $I$ in $S$, define the relation $\sim_I$ on $S$ by
\[a\sim_I b\iff\exists s\in S,\exists i,j\in I\left(a=s\lor i\text{ and }b=s\lor j\right)\]

\begin{lemma}\label{lemmaquotientinversesemigroup}
If $S$ is a $\lor$-ideal of a distributive inverse semigroup $S$, then
\begin{enumerate}
\item[(a)] If $S$ has a zero then $0\in I$;
\item[(b)] $i\in I$ if and only if $i^*\in I$;
\item[(c)] If $i\in I$ and $j\leq i$ then $j\in I$;
\item[(d)] $a\sim_I b$ if and only if $a^*\sim_I b^*$;
\item[(e)] If $a=b\lor i$ and $i\in I$ then $a\sim_I b$.
\end{enumerate}
\end{lemma}
\begin{proof}
\begin{enumerate}
\item[(a)] Letting $i\in I$ be arbitrary (because $I\neq\varnothing$), we have $0=0i\in I$.
\item[(b)] If $i\in I$ then $i^*=i^*ii^*\in I$.
\item[(c)] If $i\in I$ and $j\leq i$ then $j=ij^*j\in I$, because $I$ is absorbing.
\item[(d)] If $a\sim_I b$, choose $s\in S$ and $i,j\in I$ with $a=s\lor i$, $b=s\lor j$, so that $a^*=s^*\lor i^*$ and $b^*=s^*\lor j^*$, and thus $a^*\sim_Ib^*$ by item (b).
\item[(e)] If $a=b\lor i$ with $i\in I$, then since $b=b\lor bi^*i$ and $bi^*i\in I$ we conclude that $a\sim_I b$.\qedhere
\end{enumerate}
\end{proof}

\begin{theorem}
If $I$ is a $\lor$-ideal of a distributive inverse semigroup $S$, then $\sim_I$ is a congruence.
\end{theorem}
\begin{proof}
To see that the relation $\sim_I$ is reflexive, take $i\in I$ an arbitrary element, so that for any $a\in S$ we have $a=a\lor ai^*i$ and $ai^*i\in I$, therefore $a\sim_I a$. Moreover $\sim_I$ is immediately symmetric from its definition, so in order to check that it is an equivalence relation, suppose $a\sim_I b$ and $b\sim_I c$, so there are $s,t\in S$ and $i,j,k,l\in I$ with
\[a=s\lor i,\quad b=s\lor j=t\lor k,\quad c=t\lor l.\]
In particular, $s,t\leq b$ so $s$ and $t$ are compatible and $s\land t=st^*t$. Also, $i,sk^*k\leq a$, so $i$ and $sk^*k$ are compatible and belong to $I$ (since it is an ideal) and hence $i\lor sk^*k\in I$.

Using distributivity, we obtain
\[a=s\lor i=sb^*b\lor i=s(t^*t\lor k^*k)\lor i=st^*t\lor(sk^*k\lor i)=(s\land t)\lor(sk^*k\lor i)\]
and similarly, $l\lor tj^*j\in I$ and $c=(s\land t)\lor(l\lor tj^*j)$, which proves $a\sim_I c$ and so $\sim_I$ is an equivalence relation.

To check that $\sim_I$ is a congruence, first note that if $a\sim_I b$ and $c\in S$, then we can write
\[a=s\lor i,\quad b=s\lor j\]
for certain $s\in S$, $i,j\in I$, so by distributivity,
\[ac=sc\lor ic\quad\text{and}\quad bc=sc\lor jc\]
and therefore $ac\sim_I bc$, and similarly $ca\sim_I ba$. This fact and transitivity of $\sim_I$ imply that if $a\sim_I b$ and $c\sim_I d$, then $ac\sim_I bc\sim_I bd$, so $\sim_I$ is a congruence.\qedhere
\end{proof}

\begin{definition}\nomenclature[S]{$S/I$}{Quotient of $I$ by $S$}\index{Quotient of a distributive inverse semigroup}\nomenclature[S]{$\tildeoveri{a}$}{The class of an element $a$ in a quotient semigroup}
If $I$ is a $\lor$-ideal of a distributive inverse semigroup $S$, we denote by $S/I$ the \emph{quotient semigroup} of $S$ by $\sim_I$. The $\sim_I$-class of an element $s\in S$ will be denoted $\tildeoveri{s}$. We denote by $\pi_I:S\to S/I$ the quotient map.
\end{definition}

\begin{theorem}\index{Quotient of a distributive inverse semigroup}\label{theoremquotientmapofsemigroupisverygood}
Let $I$ be a $\lor$-ideal of a distributive inverse semigroup $S$. Then $S/I$ is as inverse semigroup and $(\tildeoveri{a})^*=(a^*)^{\sim I}$ for all $a\in S$. The quotient map $\pi_I:S\to S/I$ is a $\lor$ and a $\land$-morphism. Moreover, $S/I$ has a zero if and only if $I$ is downwards directed, in which case $I=\pi_I^{-1}(0)$.
\end{theorem}
\begin{proof}
Since $\pi_I$ is a surjective semigroup morphism, then Proposition \ref{propositionpropertiesofmorphismsofsemigroups} implies that $S/I$ is an inverse semigroup and $(\tildeoveri{a})^*=(a^*)^{\sim I}$

%Note that $S/I$ is regular, because it is a homomorphic image of $S$, and a commutative subsemigroup of $S/I$ by Lemma~\ref{lemmaidempotentsofquotient}, so Theorem~\ref{theoreminverseiffescommute} implies that $S/I$ is an inverse semigroup. The fact that $[a]=[a^*]$ follows because $[a^*]$ is one inverse of $[a]$, and thus it is the unique one.

%Suppose now that $[a],[b]\in S/I$ are compatible. This means that $[a^*b]\in E(S/I)=\pi(E(S))$, so $[a^*b]=[g]$ for some $g\in E(S)$, again by Lemma~\ref{lemmaidempotentsofquotient}. But then there are certain $e\in S$ and $i_1,i_2\in I$ with
%\[a^*b=e\lor i_1,\qquad\text{and}\qquad g=e\lor i_2\]
%in particular $e\leq g$ so $e\in E(S)$, $e\leq a^*b$ and $[e]=[a^*b]$. Similarly, there is $f\in E(S)$ with $f\leq ab^*$ and $[f]=[ab^*]$.
%
%Then the elements $fa,be\in S$ are compatible, since $e,f\in E(S)$ and
%\[(fa)^*be=a^*fbe\leq a^*be=e,\qquad\text{and}\qquad(fa)(be)^*=faeb^*\leq fab^*=f\]

Suppose $a,b\in S$ are compatible. Then $\left(a\lor b\right)^{\sim I}$ is an upper bound of both $\tildeoveri{a}$ and $\tildeoveri{b}$, so let us prove it is the smallest one. If $\tildeoveri{p}\geq\tildeoveri{a},\tildeoveri{b}$, where $p\in S$, then we can find $e\in E(S)$, $s\in S$, and $i_1,i_2\in I$ satisfying
\[pe=s\lor i_1,\quad a=s\lor i_2,\]
and in particular $s\leq p$. Similarly, we can find $t\in S$ and $j_2\in I$ with $t\leq p$ and $b=t\lor j_2$. Thus $s$ and $t$ are compatible and $s\lor t\leq p$, so
\[a\lor b=s\lor i_2\lor t\lor j_2=(s\lor t)\lor(i_2\lor j_2)\]
which proves $\left(a\lor b\right)^{\sim I}=\left(s\lor t\right)^{\sim I}\leq \tildeoveri{p}$.

If $a\land b$ exists, then $\left(a\land b\right)^{\sim I}$ is a lower bound of both $\tildeoveri{a}$ and $\tildeoveri{b}$, so to prove it is the largest one suppose $\tildeoveri{p}\leq \tildeoveri{a},\tildeoveri{b}$. There are idempotents $e,f\in E(S)$, elements $s,t\in S$ and $i_1,i_2,j_1,j_2\in I$ with
\[ae=s\lor i_1,\quad p=s\lor i_2,\quad bf=t\lor j_1,\quad p=t\lor j_2.\]
Moreover, all of $s,t,i_2,j_2$ are bounded by $p$ and are thus compatible and admit pairwise meets and joins, and combinations thereof, so we can apply Proposition \ref{corollarydistributivedistributemeets} to distribute meets over joins and obtain
\[p=(s\lor i_2)\land(t\lor j_2)=(s\land t)\lor k\]
for an appropriate $k\in I$. This implies $\tildeoveri{p}=(s\land t)^{\sim I}\leq(a\land b)^{\sim I}$.

If $I$ is downwards directed, then any two elements of $I$ are $\sim_I$-equivalent (to any common lower bound), and since $I$ is an ideal this implies that the class $\tildeoveri{i}$ of any $i\in I$ is the zero of $S/I$. This in fact proves that, in this case, $I\subseteq\pi_I^{-1}(0)$.

Conversely, suppose $\tildeoveri{o}$ is a zero in $S/I$, where we may assume $o\in E(S)$, and let $a,b\in I$. Since $(ao)^{\sim I}=\tildeoveri{o}=(bo)^{\sim I}$, there are $s,t\in S$ and $i_1,i_2,j_1,j_2\in I$ with
\[ao=s\lor i_1,\quad o=s\lor i_2,\quad bo=t\lor j_1,\quad\text{and } o=t\lor j_2.\]
In particular, $s\leq ao$ so $s\in I$. From $s,t\leq o$ we see that $s$ and $t$ are compatible, $s\land t=st^*t\in I$ and $s\land t$ is a lower bound of both $a$ and $b$, and so $I$ is downwards directed. Moreover, $o=s\lor i_2\in I$, which proves $\pi_I^{-1}(0)\subseteq I$ in this case.\qedhere
\end{proof}

\begin{example}
$S/I$ is not necessarily a distributive inverse semigroup. Let $X=\left\{1,2,3,4,5,6,7\right\}$, $a,b\in\mathcal{I}(X)$ be the partial bijections of $X$ defined as
\[a(1)=2,\quad a(4)=5,\qquad b(1)=3,\quad b(6)=7\]
\[\begin{tikzcd}[row sep=small]
&2&4\arrow[r,"a"]&5\\
1\arrow[ur,"a"]\arrow[dr,"b"]&&&\\
&3&6\arrow[r,"b"]&7
\end{tikzcd}\]
and then consider the subsemigroup $S\subseteq\mathcal{I}(X)$ of partial bijections $f:\dom(f)\to\ran(f)$ such that
\begin{itemize}
\item if $4\in\operatorname{dom}(f)$ then $a\leq f$ or $a^*a\leq f$;
\item if $5\in\operatorname{dom}(f)$ then $a^*\leq f$ or $aa^*\leq f$;
\item if $6\in\operatorname{dom}(f)$ then $b\leq f$ or $b^*b\leq f$;
\item if $7\in\operatorname{dom}(f)$ then $b^*\leq f$ or $bb^*\leq f$.
\end{itemize}
Let us prove that this is a sub-inverse semigroup of $\mathcal{I}(X)$: Suppose $f,g\in S$, and that $4\in\operatorname{dom}(fg)$. Then $4\in\operatorname{dom}(g)$. There are two possibilities:
\begin{itemize}
\item If $a\leq g$, then $g(4)=a(4)=5\in\operatorname{dom}(f)$, and there are two sub-possibilities:
\begin{itemize}
\item if $a^*\leq f$, then $a^*a\leq fg$;
\item if $aa^*\leq f$, then $a=aa^*a\leq fg$;
\end{itemize}
\item If $a^*a\leq g$, then $g(4)=4\in\operatorname{dom}(f)$, and there are two sub-possibilities:
\begin{itemize}
\item If $a\leq f$, then $a=aa^*a\leq fg$;
\item If $a^*a\leq f$, then $a^*a=a^*aa^*a\leq fg$;
\end{itemize}
\end{itemize}
In any case, either $a\leq fg$ or $a^*a\leq fg$ whenever $4\in\operatorname{dom}(fg)$. The other cases are dealt with similarly. Also, note that $S$ is closed under inverses of $\mathcal{I}(X)$.

Now we can prove that $S$ is closed under joins of $\mathcal{I}(X)$. Suppose that $f,g\in S$ are compatible, so the join $f\lor g$ exists in $\mathcal{I}(X)$. Again, suppose that $4\in\operatorname{dom}(f\lor g)$, so without loss of generality let us assume further that $4\in\operatorname{dom}(f)$. Then either $a\leq f\leq f\lor g$ or $a^*a\leq f\leq f\lor g$. The other cases are dealt with similarly.

Let $I=\left\{f\in S:\operatorname{dom}(f)\subseteq\left\{1,2,3\right\}\right\}$. The restriction map $\theta:S\to\mathcal{I}(\left\{4,5,6,7\right\})$, $\theta(f)=f|_{\operatorname{dom}(f)\cap\left\{4,5,6,7\right\}}$, is a morphism and $\theta(f)=\theta(g)\iff f\sim_I g$, so $\theta$ factors through an injective morphism of $S/I$ into $\mathcal{I}(\left\{4,5,6,7\right\})$. 

However $\theta(a)$ and $\theta(b)$ are compatible but do not have a common join in $\theta(S)$. Indeed, suppose that $\theta(f)=\theta(a)\lor\theta(b)$, where $f\in S$. Then $f(4)=a(4)=5$, which implies, by the definition of $S$, that $a\leq f$. However similarly, $f(6)=b(6)=7$ which implies $b\leq f$, but $a$ and $b$ are not compatible in $S$, a contradiction.

This proves that $\theta(S)$ - and therefore $S/I$ - is not a distributive inverse semigroup.
\end{example}

On the other hand, if $S$ is (weakly) Boolean then $S/I$ is also (weakly) Boolean. In the previous example, note that the restriction $a|_{\left\{1\right\}}$ does not admit a relative complement in $a$.

\begin{theorem}\label{theoremcompatibleinquotientareimageofcompatible}
If $S$ is a weakly Boolean inverse semigroup and $I$ is a $\lor$-ideal of $S$, then for all pair of compatible elements $\alpha,\beta\in S/I$ there are compatible elements $a,b\in S$ such that $\tildeoveri{a}=\alpha$ and $\tildeoveri{b}=\beta$.
\end{theorem}
\begin{proof}
Suppose $\alpha=\tildeoveri{a}$ and $\beta=\tildeoveri{b}$ are compatible in $S/I$. By Proposition \ref{propositionpropertiesofmorphismsofsemigroups}(e), we can find $e\in E(S)$ and $i\in I$ such that $b^*a=e\lor i$. Note that $bi\leq bb^*a\leq a$, so using Proposition \ref{propositionrelativecomplementofsmallerelementinweaklybooleansemigroup}, we can consider the element $a_0=a\setminus bi$.

Similarly, there are $f\in E(S)$ and $j\in I$ such that $ba^*=f\lor j$, and we can consider the element $b_0=b\setminus ja$. Let us prove that $a_0$ and $b_0$ are compatible:
\[b_0^*a_0\leq b^*(a\setminus bi)=(b^*a)\setminus(b^*b i)\]
but since $i\leq b^*a$ then $b^*bi=i$. Applying Proposition \ref{propositionrelativecomplementsandoperations}(c),
\[b_0^*a_0\leq b^*a\setminus i\leq e\]
so $b_0^*a_0\in E(S)$. Similarly, $b_0 a_0^*\in E(S)$, so $a_0$ and $b_0$ are compatible. Moreover, $a=a_0\lor bi$ and $b=b_0\lor ja$, so from \ref{lemmaquotientinversesemigroup}(e) we have
\[\alpha=\tildeoveri{a}=\tildeoveri{a_0}\qquad\text{and}\qquad\beta=\tildeoveri{b}=\tildeoveri{b_0}\qedhere\]
\end{proof}
An immediate consequence of this and Theorem \ref{theoremquotientmapofsemigroupisverygood} follows.

\begin{theorem}\label{theoremquotientofbooleanisboolean}
If $S$ is a weakly Boolean inverse semigroup and $I$ is a $\lor$-ideal of $S$, then $S/I$ is weakly Boolean, and the quotient map $\pi_I:S\to S/I$ is a $\lor$-morphism. If $S$ is Boolean then $S/I$ is Boolean and the quotient map $\pi_I:S\to S/I$ is a morphism of Boolean inverse semigroups.
\end{theorem}

\begin{proposition}\label{propositionquotientrelationofbooleansemigroup}
If $S$ is a Boolean inverse semigroup and $I\subseteq S$ is a $\lor$-ideal, then $a\sim_I b$ if and only if $a\setminus b,b\setminus a\in I$.
\end{proposition}
\begin{proof}
On one hand, if $a\sim_I b$ then take $s\in S$ and $i,j\in I$ with $a=s\lor i$, $b=s\lor j$. Then by Proposition \ref{propositionrelativecomplementsandoperations}(c) and (d),
\[a\setminus b\leq a\setminus s\leq i\]
so $a\setminus b\in I$ by Lemma \ref{lemmaquotientinversesemigroup}(c), and similarly $b\setminus a\in I$.

Conversely, if $a\setminus b,b\setminus a\in I$ then $a=(a\land b)\lor(a\setminus b)$ and $b=(a\land b)\lor(b\setminus a)$ implies $a\sim_I b$.\qedhere
\end{proof}

To finish this section, we describe $\lor$-ideals of distributive inverse semigroups as the zero sets of semigroup morphisms, and prove an Isomorphism Theorem for $\lor$-ideals of Boolean inverse semigroups. The simpler direction of the following proposition was already noticed in \cite[Proposition 2.10]{MR3519501}, without proof.

\begin{proposition}
A nonempty subset $I$ of a distributive inverse semigroup $S$ with zero is a $\lor$-ideal if and only if there is an inverse semigroup $T$ with $0$ and a $\lor$-morphism $\theta:S\to T$ with $I=\theta^{-1}(0)$.
\end{proposition}
\begin{proof}
It is easy enough to verify that $\theta^{-1}(0)$ is a $\lor$-ideal, under the assumptions above, and the converse direction follows using the canonical morphism $S\to S/I$ and Theorem \ref{theoremquotientmapofsemigroupisverygood}.\qedhere
\end{proof}

\begin{theorem}[Isomorphism Theorem for Boolean inverse semigroups]\index{{Isomorphism Theorem for Boolean inverse semigroups}}
Let $\theta:S\to T$ be a morphism of Boolean inverse semigroups (as in Definition \ref{definitionmorphismofbooleaninversesemigroups}) and $I\subseteq\theta^{-1}(0)$ a $\lor$-ideal. Then
\begin{enumerate}
\item[(a)] $\theta$ factors uniquely to a Boolean inverse semigroup morphism $\theta_I:S/I\to T$;
\item[(b)] $I=\theta^{-1}(0)$ if and only if $\theta_I$ is injective.
\end{enumerate}
\end{theorem}
\begin{proof}
First note that if $I\subseteq\theta^{-1}(0)$ then $\theta(0)=0$ because $0\in I$.

\begin{enumerate}
\item[(a)] We just need to prove that $a\sim_I b$ implies $\theta(a)=\theta(b)$. If $a\sim_I b$ then choose $s\in S$ and $i,j\in I$ with
\[a=s\lor i\qquad\text{and}\qquad b=s\lor j\]
so
\[\theta(a)=\theta(s)\lor\theta(i)=\theta(s)\lor 0=\theta(s)\lor\theta(j)=\theta(b)\]
so $\theta$ factors through a map $\theta_I:S/I\to T$, and the factor map is unique because the quotient map $S\to S/I$ is surjective. To see that $\theta_I$ is a morphism of Boolean inverse semigroups, apply Theorem \ref{theoremquotientmapofsemigroupisverygood} and the fact that compatible elements in $S/I$ are images of compatible elements in $S$ (by Theorem \ref{theoremcompatibleinquotientareimageofcompatible}).
\item[(b)] Let us prove that if $I=\theta^{-1}(0)$ then $\theta_I$ is injective. If $\theta(a)=\theta(b)$, then
\[\theta(a\setminus b)=\theta(a)\setminus\theta(b)=\theta(a)\setminus\theta(a)=0\]
so $a\setminus b\in\theta^{-1}(0)=I$, and similarly $b\setminus a\in I$, which means by \ref{propositionquotientrelationofbooleansemigroup} that $a\sim_I b$. This means that $\theta_I$ is injective.

Conversely, if $\theta_I$ is injective and $a\in\theta^{-1}(0)$ then
\[\theta_I(\tildeoveri{a})=\theta(a)=0=\theta_I(\tildeoveri{0})\]
so $a\sim_I 0$, which means that $a\in I$.\qedhere
\end{enumerate}
\end{proof}